% =====================================================================
% UTFPR - Ciências do Ambiente
% Projeto 1: Mobilidade Urbana e Sustentabilidade no Centro de Curitiba
% Formato ABNT
% =====================================================================
\documentclass[12pt,a4paper,oneside]{abntex2}

% --- Pacotes ---
\usepackage[utf8]{inputenc}
\usepackage[T1]{fontenc}
\usepackage[brazil]{babel}
\usepackage{amsmath,amssymb}
\usepackage{graphicx}
\usepackage{float}
\usepackage{indentfirst}
\usepackage{setspace}
\usepackage{hyperref}
\usepackage{enumitem}
\usepackage{tabularx}
\usepackage{booktabs}
\usepackage{longtable}
\usepackage{array}
\usepackage{multirow}
\usepackage{pdflscape}

% --- Metadados ---
\titulo{Projeto de Mobilidade Urbana e Sustentabilidade no Centro de Curitiba}
\tituloestrangeiro{Urban Mobility and Sustainability Project in Downtown Curitiba}
\autor{Equipe de Engenharia Multidisciplinar}
\local{Curitiba}
\data{2026}
\instituicao{Universidade Tecnológica Federal do Paraná -- UTFPR}
\tipotrabalho{Relatório Técnico de Projeto}
\preambulo{Projeto interdisciplinar apresentado à Unidade Curricular de Ciências do Ambiente do curso de Engenharia da Universidade Tecnológica Federal do Paraná (UTFPR), como requisito parcial para aprovação na disciplina.}

% --- Configurações ---
\OnehalfSpacing
\setlength{\parindent}{1.25cm}

% --- Hyperref ---
\hypersetup{
  colorlinks=true,
  linkcolor=black,
  citecolor=black,
  urlcolor=blue
}

\begin{document}

% =====================================================================
% ELEMENTOS PRÉ-TEXTUAIS
% =====================================================================

% --- Capa ---
\imprimircapa

% --- Folha de rosto ---
\imprimirfolhaderosto

% --- Ficha catalográfica (opcional, omitida) ---

% --- Folha de aprovação (omitida - uso interno) ---

% --- Resumo ---
\begin{resumoumacoluna}
Este projeto tem como objetivo diagnosticar e mitigar os gargalos de mobilidade urbana no quadrilátero central de Curitiba, definido pelo perímetro entre as ruas Presidente Kennedy, Calçadão da XV, Marechal Floriano Peixoto, Martim Afonso, Mariano Torres e Mário Tourinho. A intervenção proposta integra soluções das engenharias de Computação, Elétrica e Mecânica, utilizando modelagem computacional de tráfego, infraestrutura de sensores IoT com controle semafórico adaptativo (Onda Verde) e análise de eficiência energética para quantificar a redução de emissões de CO\textsubscript{2} e material particulado. O trabalho segue quatro fases: levantamento de campo, modelagem e simulação, proposta de intervenção e validação dos resultados. Espera-se comprovar, por simulação em software SUMO, redução de pelo menos 15\% no tempo de deslocamento durante os horários de pico e diminuição proporcional nas emissões veiculares.

\textbf{Palavras-chave:} mobilidade urbana; sustentabilidade; semáforos inteligentes; IoT; emissões veiculares; Curitiba.
\end{resumoumacoluna}

% --- Listas ---
\pdfbookmark[0]{\listfigurename}{lof}
\listoffigures*
\cleardoublepage

\pdfbookmark[0]{\listtablename}{lot}
\listoftables*
\cleardoublepage

\pdfbookmark[0]{\contentsname}{toc}
\tableofcontents*
\cleardoublepage

% =====================================================================
% ELEMENTOS TEXTUAIS
% =====================================================================

% --- 1. INTRODUÇÃO ---
\chapter{Introdução}

O crescimento urbano desordenado associado à predominância do transporte individual motorizado configura-se como um dos maiores desafios das metrópoles contemporâneas. Em Curitiba, cidade reconhecida internacionalmente por seu planejamento urbano inovador, o quadrilátero central delimitado pelas ruas Presidente Kennedy, Calçadão da XV de Novembro, Marechal Floriano Peixoto, Martim Afonso, Mariano Torres e Mário Tourinho concentra intensa atividade comercial, institucional e de serviços, resultando em congestionamentos críticos nos horários de pico --- das 06h30 às 09h00 e das 17h00 às 19h00.

A elevada densidade veicular nessa região não apenas compromete a fluidez do tráfego, como também intensifica a emissão de gases poluentes, o consumo de combustíveis fósseis e os níveis de ruído urbano, impactando negativamente a qualidade de vida da população e o meio ambiente. Diante desse cenário, o presente projeto propõe uma intervenção técnica multidisciplinar que integra as competências das engenharias de Computação, Elétrica e Automação, e Mecânica para conceber, simular e validar soluções de mobilidade sustentável aplicáveis à realidade curitibana.

A abordagem adotada fundamenta-se na premissa de que a engenharia não se limita à construção de máquinas ou \textit{softwares}, mas ao desenho de cidades onde as pessoas possam viver e respirar melhor. Para tanto, este trabalho estrutura-se em quatro capítulos principais, além desta introdução: o Capítulo 2 apresenta a justificativa e os objetivos do projeto; o Capítulo 3 desenvolve a fundamentação teórica; o Capítulo 4 descreve a metodologia estruturada em quatro fases; o Capítulo 5 detalha as soluções técnicas por especialidade; o Capítulo 6 apresenta os resultados esperados e, por fim, o Capítulo 7 sintetiza as conclusões.

% --- 2. JUSTIFICATIVA E OBJETIVOS ---
\chapter{Justificativa e Objetivos}

\section{Justificativa}

A escolha do quadrilátero central de Curitiba como objeto de estudo justifica-se por três fatores convergentes. Primeiro, trata-se de uma região classificada como Setor Especial Estrutural pelo Instituto de Pesquisa e Planejamento Urbano de Curitiba (IPPUC), o que lhe confere importância estratégica para a mobilidade metropolitana. Segundo, os dados da Urbanização de Curitiba S.A. (URBS) indicam volumes diários de tráfego superiores a 50 mil veículos nas vias principais do perímetro, com redução da velocidade média para menos de 12 km/h nos horários de pico. Terceiro, a presença do Calçadão da XV --- eixo exclusivo para pedestres --- gera conflitos adicionais entre fluxos veiculares e travessias de pedestres, agravando a retenção nas interseções periféricas.

Do ponto de vista ambiental, estima-se que cada veículo leve em ciclo ``para-e-anda'' consuma até 40\% mais combustível do que em regime de cruzeiro, elevando proporcionalmente as emissões de dióxido de carbono (CO\textsubscript{2}), monóxido de carbono (CO), óxidos de nitrogênio (NO\textsubscript{x}) e material particulado (MP\textsubscript{2,5} e MP\textsubscript{10}). A Organização Mundial da Saúde (OMS) recomenda que a concentração média anual de MP\textsubscript{2,5} não ultrapasse 5 $\mu$g/m\textsuperscript{3}, valor frequentemente excedido em centros urbanos congestionados.

\section{Objetivo Geral}

Propor e validar tecnicamente um conjunto de intervenções de engenharia de tráfego capazes de reduzir o tempo de deslocamento veicular e as emissões atmosféricas no quadrilátero central de Curitiba durante as janelas de pico, por meio da integração de sistemas inteligentes de transporte (ITS), simulação computacional e análise de eficiência energética.

\section{Objetivos Específicos}

\begin{enumerate}[label=\alph*)]
  \item Realizar o levantamento do Volume Plano Diário (VPD) nas vias principais do perímetro e identificar os pontos críticos de retenção;
  \item Construir um modelo computacional do cenário atual (As-Is) utilizando o \textit{software} SUMO (Simulation of Urban Mobility);
  \item Projetar um sistema de controle semafórico adaptativo baseado em sensores IoT e lógica de Onda Verde;
  \item Modelar matematicamente o perfil de emissões de CO\textsubscript{2} da frota circulante em função dos ciclos de aceleração e frenagem;
  \item Simular o cenário proposto (To-Be) e quantificar os ganhos percentuais de fluidez e redução de emissões.
\end{enumerate}

% --- 3. FUNDAMENTAÇÃO TEÓRICA ---
\chapter{Fundamentação Teórica}

\section{Mobilidade Urbana e Sistemas Inteligentes de Transporte}

Os Sistemas Inteligentes de Transporte (ITS) constituem um campo interdisciplinar que aplica tecnologias de informação, comunicação e sensoriamento à gestão da mobilidade urbana. Segundo \cite{IBGE2022}, o Brasil possuía, em 2022, frota superior a 115 milhões de veículos, dos quais 59\% eram automóveis de passeio, evidenciando a urgência de soluções tecnológicas para mitigar congestionamentos.

A modelagem computacional de tráfego apoia-se em três paradigmas principais: modelos macroscópicos, que tratam o fluxo veicular como fluido contínuo; modelos microscópicos, que simulam o comportamento individual de cada veículo; e modelos mesoscópicos, que operam em nível intermediário. O SUMO (\textit{Simulation of Urban Mobility}), desenvolvido pelo \textit{German Aerospace Center} (DLR), é uma plataforma \textit{open-source} amplamente adotada em pesquisas acadêmicas por sua capacidade de simulação microscópica, suporte a semáforos adaptativos e integração com dados reais de sensoriamento \cite{SUMO2025}.

\section{Semáforos Adaptativos e Onda Verde}

O conceito de Onda Verde consiste na sincronização progressiva de semáforos ao longo de um eixo viário, de modo que um veículo trafegando à velocidade de projeto encontre sinais verdes consecutivos, eliminando paradas desnecessárias. Matematicamente, a defasagem temporal $\Delta t$ entre dois semáforos consecutivos separados por distância $d$ é dada pela Equação \ref{eq:onda_verde}:

\begin{equation}
\Delta t = \frac{d}{v_p}
\label{eq:onda_verde}
\end{equation}

onde $v_p$ é a velocidade de progressão desejada (tipicamente entre 40 e 50 km/h em vias arteriais urbanas).

Os semáforos adaptativos vão além da sincronização fixa, empregando sensores indutivos, câmeras ou laços magnéticos para ajustar dinamicamente os tempos de verde e vermelho com base no fluxo real de veículos. O algoritmo de controle em malha fechada mais difundido para essa finalidade é o controle PID (\textit{Proportional-Integral-Derivative}), cuja equação é:

\begin{equation}
u(t) = K_p e(t) + K_i \int_0^t e(\tau) d\tau + K_d \frac{de(t)}{dt}
\label{eq:pid}
\end{equation}

onde $u(t)$ é o sinal de controle (duração do ciclo semafórico), $e(t)$ é o erro entre o fluxo desejado e o medido, e $K_p$, $K_i$, $K_d$ são os ganhos do controlador.

\section{Emissões Veiculares e Eficiência Energética}

O perfil de emissões de um veículo em ambiente urbano é fortemente influenciado pelo regime de condução. O ciclo ``para-e-anda'' (\textit{stop-and-go}) multiplica os eventos de aceleração, nos quais o motor opera fora de seu ponto ótimo de eficiência. A modelagem das emissões pode ser realizada pelo modelo COPERT (\textit{Computer Programme to Calculate Emissions from Road Transport}), adotado pela União Europeia, ou pela abordagem baseada em potência específica veicular (VSP -- \textit{Vehicle Specific Power}), conforme a Equação \ref{eq:vsp}:

\begin{equation}
\text{VSP} = v \cdot (a \cdot (1 + \varepsilon) + g \cdot \text{sen}(\theta) + g \cdot C_r) + \frac{1}{2} \rho_a \cdot C_D \cdot A_f \cdot v^3 / m
\label{eq:vsp}
\end{equation}

onde $v$ é a velocidade instantânea, $a$ a aceleração, $\varepsilon$ o fator de massa rotativa, $g$ a aceleração gravitacional, $\theta$ a inclinação da via, $C_r$ o coeficiente de resistência ao rolamento, $\rho_a$ a densidade do ar, $C_D$ o coeficiente de arrasto, $A_f$ a área frontal e $m$ a massa do veículo.

A norma ABNT NBR 6601:2021 estabelece os limites de emissão para veículos leves no Brasil, enquanto a NBR 10151:2019 regula os níveis de ruído em ambientes externos, fixando em 55 dB(A) o limite diurno para áreas mistas com vocação comercial e administrativa \cite{ABNT10151}.

\section{Internet das Coisas (IoT) no Monitoramento de Tráfego}

A Internet das Coisas (IoT) viabiliza a coleta descentralizada de dados de tráfego por meio de dispositivos de baixo custo e alta capilaridade. Sensores infravermelhos, laços indutivos embutidos no pavimento, câmeras com visão computacional e detectores de \textit{bluetooth}/Wi-Fi compõem um ecossistema de monitoramento que alimenta a malha de controle semafórica com dados em tempo real. Protocolos como MQTT (\textit{Message Queuing Telemetry Transport}) e CoAP (\textit{Constrained Application Protocol}) são particularmente adequados à transmissão de telemetria veicular em redes de baixa largura de banda.

% --- 4. METODOLOGIA ---
\chapter{Metodologia}

A metodologia deste projeto estrutura-se em quatro fases sequenciais e interdependentes, conforme o \textit{roadmap} estabelecido no manual do projeto. Cada fase é descrita a seguir com seus respectivos procedimentos, ferramentas e entregáveis.

\section{Fase 1 -- Levantamento de Campo e Diagnóstico}

\subsection{Procedimentos}

\begin{enumerate}[label=\textbf{P\arabic*:},leftmargin=2cm]
  \item \textbf{Contagem volumétrica classificada:} utilizando contadores manuais e câmeras com reconhecimento de placas (LPR -- \textit{License Plate Recognition}), realizou-se a contagem de veículos nas vias principais do quadrilátero durante as janelas de pico matutina (06h30--09h00) e vespertina (17h00--19h00). Os veículos foram classificados em quatro categorias: motocicletas, automóveis, ônibus e caminhões.
  \item \textbf{Mapeamento geométrico:} levantamento das larguras de faixa, raios de curvatura nas interseções, comprimento das baias de carga/descarga, e distância entre semáforos consecutivos.
  \item \textbf{Identificação de gargalos:} registro fotográfico e georreferenciado dos pontos onde o fluxo é obstruído por estacionamento irregular, operações de carga/descarga fora do horário permitido ou conflitos com pedestres nas proximidades do Calçadão da XV.
  \item \textbf{Medição de ruído ambiental:} utilizando decibelímetros classe 2 (IEC 61672), foram realizadas medições nos pontos de maior retenção, registrando os níveis de pressão sonora equivalente ($L_{Aeq}$) em intervalos de 15 minutos.
\end{enumerate}

\subsection{Ferramentas}

Os dados coletados foram consolidados em planilhas eletrônicas e exportados para o formato compatível com o SUMO (arquivos \texttt{.net.xml} para a rede e \texttt{.rou.xml} para as rotas). Os softwares utilizados nesta fase incluem o QGIS para geoprocessamento e o \textit{OpenStreetMap} como base cartográfica.

\subsection{Entregável}

Relatório de diagnóstico contendo: mapa de calor do VPD por trecho, tabela de tempos de ciclo semafórico atuais, planta de obstáculos geométricos e mapa de ruído ($L_{Aeq}$).

\section{Fase 2 -- Modelagem e Simulação (Cenário As-Is)}

\subsection{Procedimentos}

\begin{enumerate}[label=\textbf{P\arabic*:},leftmargin=2cm]
  \item \textbf{Construção da rede viária:} a malha viária do quadrilátero foi modelada no SUMO a partir de dados do \textit{OpenStreetMap}, refinados manualmente com as medições de campo. Cada interseção foi parametrizada com o número de faixas, sentidos permitidos e tempos semafóricos atuais.
  \item \textbf{Geração da demanda:} com base nos dados de VPD, foram geradas matrizes origem-destino (OD) para cada janela de pico, calibradas para que o volume simulado de veículos em cada trecho correspondesse ao observado em campo (erro máximo admissível de $\pm 10\%$).
  \item \textbf{Definição da frota circulante:} a composição da frota foi parametrizada conforme dados do IPPUC e da URBS, com distribuição de categorias veiculares e fatores de emissão específicos para cada classe.
  \item \textbf{Execução da simulação:} o cenário base foi executado no SUMO com passo de simulação de 0,1 segundo, gerando métricas de tempo de viagem, velocidade média, comprimento de fila e emissões (por meio do módulo \texttt{emissionsDrivingCycle}).
\end{enumerate}

\subsection{Ferramentas}

SUMO (v1.20 ou superior), Python 3.11+ (com bibliotecas TraCI para interface com o SUMO e pandas para tratamento de dados).

\subsection{Entregável}

Arquivos de simulação do cenário \textit{As-Is} e tabela-resumo das métricas de desempenho e emissões para cada via do perímetro.

\section{Fase 3 -- Proposta de Intervenção}

\subsection{Subproposta A -- Engenharia de Computação}

Desenvolvimento de um \textit{framework} de otimização para semáforos adaptativos. O sistema emprega lógica \textit{fuzzy} para tomada de decisão, pois o tráfego urbano apresenta características não lineares que o controle PID tradicional não captura adequadamente. As variáveis linguísticas de entrada são: comprimento de fila ($QF$), taxa de ocupação da via ($O$) e diferença de fluxo entre ciclos ($\Delta F$). As regras de inferência foram elaboradas a partir do conhecimento de especialistas da URBS, gerando um sinal de saída que determina a extensão ou redução do tempo de verde na próxima fase.

Ademais, foi implementado um módulo de predição de curto prazo utilizando redes LSTM (\textit{Long Short-Term Memory}), treinadas com séries temporais de fluxo de 30 dias, capazes de antecipar picos de demanda com horizonte de 5 minutos.

\subsection{Subproposta B -- Engenharia Elétrica e Automação}

O projeto de infraestrutura prevê a implantação de:

\begin{itemize}
  \item \textbf{Laços indutivos:} 36 pontos de detecção (um por faixa de rolamento em cada aproximação de interseção), operando a 20--120 kHz, com controlador lógico programável (CLP) local responsável pelo pré-processamento do sinal e comunicação com o nó concentrador via RS-485.
  \item \textbf{Câmeras IP:} 12 câmeras com sensor CMOS de 4 MP, lente varifocal de 2,8--12 mm e capacidade de visão noturna, posicionadas nas interseções de maior complexidade (Pres.~Kennedy $\times$ Mariano Torres, Mal.~Floriano $\times$ Martim Afonso e XV de Novembro $\times$ Mário Tourinho). As imagens são processadas por um algoritmo de \textit{deep learning} (YOLOv8) para contagem e classificação veicular em tempo real.
  \item \textbf{Nó concentrador:} \textit{gateway} IoT baseado em Raspberry Pi 5 com conectividade 4G/5G, executando o protocolo MQTT para publicação dos dados de sensoriamento em um \textit{broker} central.
  \item \textbf{Sincronização Onda Verde:} o controlador mestre, instalado no centro de controle da URBS, recebe os dados agregados de todos os nós e calcula, a cada 3 minutos, a defasagem ótima $\Delta t_i$ para cada par de interseções consecutivas. O sinal de sincronismo é transmitido via fibra óptica apagada da Copel Telecom para os controladores locais.
\end{itemize}

\subsection{Subproposta C -- Engenharia Mecânica}

A análise de eficiência energética concentrou-se em dois eixos:

\begin{enumerate}[label=\textbf{E\arabic*:},leftmargin=1.5cm]
  \item \textbf{Modelagem de emissões:} a partir do perfil de velocidade obtido na simulação SUMO (cenário As-Is), foram calculadas as emissões de CO\textsubscript{2}, CO, NO\textsubscript{x} e HC (hidrocarbonetos) para cada veículo em cada trecho, utilizando o modelo HBEFA (\textit{Handbook Emission Factors for Road Transport}) na versão 4.2, que fornece fatores de emissão em g/km estratificados por categoria veicular, norma de emissão (PROCONVE) e regime de tráfego.
  \item \textbf{Cálculo da economia de combustível:} o consumo total de combustível $C_{total}$ em litros para o período de simulação é dado por:
\end{enumerate}

\begin{equation}
C_{total} = \sum_{i=1}^{N} \sum_{j=1}^{T_i} \frac{\text{VSP}_{ij} \cdot \eta_{ij}}{\rho_{comb} \cdot \text{PCI}}
\label{eq:consumo}
\end{equation}

onde $N$ é o número de veículos na simulação, $T_i$ é o número de passos de tempo do veículo $i$, $\eta_{ij}$ é a eficiência instantânea do motor, $\rho_{comb}$ é a densidade do combustível e PCI é o poder calorífico inferior.

\subsection{Entregável}

Projeto executivo contendo: memória de cálculo das emissões, diagramas elétricos unifilares, arquitetura de \textit{software}, tabela-verdade da lógica \textit{fuzzy} e relação de equipamentos com especificações técnicas.

\section{Fase 4 -- Validação dos Resultados (Cenário To-Be)}

\subsection{Procedimentos}

\begin{enumerate}[label=\textbf{P\arabic*:},leftmargin=2cm]
  \item \textbf{Reconfiguração da rede no SUMO:} os tempos semafóricos fixos foram substituídos pelo controlador adaptativo desenvolvido na Fase 3, implementado via script Python/TraCI que se comunica com o SUMO a cada ciclo de simulação.
  \item \textbf{Execução do cenário To-Be:} a mesma demanda veicular do cenário As-Is foi submetida ao novo regime de controle semafórico, garantindo comparabilidade direta.
  \item \textbf{Cálculo dos indicadores de desempenho:} para cada via principal, foram computados os seguintes indicadores:

  \begin{itemize}
    \item $\Delta T$ = redução percentual do tempo médio de viagem;
    \item $\Delta V$ = aumento percentual da velocidade média;
    \item $\Delta E_{CO_2}$ = redução percentual das emissões de CO\textsubscript{2};
    \item $\Delta R$ = redução do nível de ruído equivalente ($L_{Aeq}$).
  \end{itemize}

  \item \textbf{Validação estatística:} aplicou-se o teste $t$ de Student pareado (nível de significância $\alpha = 0,05$) para verificar se as diferenças observadas eram estatisticamente significantes.
\end{enumerate}

\subsection{Entregável}

Relatório comparativo As-Is $\times$ To-Be com gráficos de barras para cada indicador e tabela de significância estatística.

% --- 5. SOLUÇÕES POR ESPECIALIDADE ---
\chapter{Plano de Projeto e Soluções por Especialidade}

Este capítulo consolida o plano de projeto na forma de cronograma e aprofunda as soluções técnicas desenvolvidas por cada especialidade da engenharia.

\section{Cronograma de Execução (Plano de Projeto)}

O projeto foi planejado para execução ao longo do semestre letivo 2026/1, com duração total de 18 semanas. A Tabela \ref{tab:cronograma} apresenta a distribuição temporal das atividades.

\begin{table}[H]
\centering
\caption{Cronograma de execução do projeto}
\label{tab:cronograma}
\small
\begin{tabularx}{\textwidth}{p{5cm} X c c c c}
\toprule
\textbf{Atividade} & \textbf{Descrição} & \textbf{Semanas} & \textbf{Início} & \textbf{Término} \\
\midrule
Kick-off e definição de papéis & Formação dos grupos e designação por especialidade & 1--2 & Semana 1 & Semana 2 \\
Levantamento de campo & Coleta de VPD, ruído e mapeamento geométrico & 3--5 & Semana 3 & Semana 5 \\
Consolidação dos dados & Tratamento estatístico e exportação para SUMO & 6--7 & Semana 6 & Semana 7 \\
Modelagem As-Is & Simulação do cenário atual & 8--10 & Semana 8 & Semana 10 \\
Projeto das soluções & Desenv. de semáforo adaptativo, IoT, emissões & 10--13 & Semana 10 & Semana 13 \\
Simulação To-Be & Validação do novo cenário & 14--15 & Semana 14 & Semana 15 \\
Análise comparativa & Cálculo de indicadores e teste estatístico & 16 & Semana 16 & Semana 16 \\
Redação do relatório & Elaboração do documento final & 16--17 & Semana 16 & Semana 17 \\
Apresentação oral & Defesa do projeto perante banca & 18 & Semana 18 & Semana 18 \\
\bottomrule
\end{tabularx}
\end{table}

\section{Solução de Engenharia de Computação}

\subsection{Modelagem Matemática do Tráfego}

O modelo microscópico adotado para a simulação no SUMO é o modelo de \textit{car-following} de Krauß, cuja equação fundamental é:

\begin{equation}
v_{seg}(t) = v_l(t) + \frac{g(t) - v_l(t) \cdot T_r}{\frac{v_l(t) + v_f(t)}{2b} + T_r}
\label{eq:krauss}
\end{equation}

onde $v_{seg}$ é a velocidade segura para o veículo seguidor, $v_l$ é a velocidade do veículo líder, $v_f$ é a velocidade do veículo seguidor, $g$ é o \textit{gap} entre os veículos, $T_r$ é o tempo de reação do motorista (calibrado em 1,0 s para condições urbanas), e $b$ é a desaceleração máxima (4,5 m/s\textsuperscript{2}).

\subsection{Algoritmo de Otimização Semafórica}

O \textit{framework} de otimização semafórica desenvolvido opera em três camadas hierárquicas, conforme ilustrado na Figura \ref{fig:arquitetura}:

\begin{figure}[H]
\centering
\fbox{\begin{minipage}{0.88\textwidth}
\vspace{0.5cm}
\begin{center}
\textbf{Arquitetura em Três Camadas do Sistema ITS}
\vspace{0.3cm}

\textbf{Camada 1 -- Aquisição:} Sensores IoT (laços, câmeras) $\rightarrow$ MQTT $\rightarrow$ \textit{Data Lake}
\vspace{0.2cm}

\textbf{Camada 2 -- Processamento:} \textit{Streaming} (Apache Kafka) $\rightarrow$ Lógica \textit{fuzzy} + LSTM $\rightarrow$ Decisão de ciclo
\vspace{0.2cm}

\textbf{Camada 3 -- Atuação:} Controladores locais $\rightarrow$ Grupos focais semafóricos $\rightarrow$ Onda Verde
\vspace{0.3cm}
\end{center}
\end{minipage}}
\caption{Diagrama de blocos da arquitetura do sistema ITS proposto}
\label{fig:arquitetura}
\end{figure}

A lógica \textit{fuzzy} descrita na Fase 3 foi implementada em Python utilizando a biblioteca \texttt{scikit-fuzzy}, com 27 regras de inferência e funções de pertinência triangulares e trapezoidais. A base de conhecimento foi refinada iterativamente comparando as decisões do sistema com as de um operador humano experiente em 100 cenários de tráfego distintos.

\subsection{Análise de Big Data}

Os dados brutos de sensoriamento (estimados em aproximadamente 2,4 GB/dia considerando taxa de 1 Hz por sensor e 36 sensores) são armazenados em um \textit{data lake} baseado em Apache Parquet e analisados com Apache Spark para gerar relatórios diários de fluxo e identificar padrões sazonais (dias úteis $\times$ fins de semana, período letivo $\times$ férias escolares).

\section{Solução de Engenharia Elétrica e Automação}

\subsection{Topologia da Rede de Sensores}

A rede de sensores IoT adota topologia estrela com \textit{backbone} em fibra óptica, conforme esquema da Figura \ref{fig:topologia}. Os laços indutivos são conectados a CLPs Siemens S7-1200, que executam a lógica de pré-processamento e enviam os pacotes de dados ao nó concentrador via Modbus TCP. As câmeras IP (Hikvision DS-2CD2T46G2) executam o modelo YOLOv8 embarcado em um \textit{edge computer} Jetson Orin Nano, enviando apenas metadados (contagem e classificação) para o \textit{broker} MQTT.

\begin{figure}[H]
\centering
\fbox{\begin{minipage}{0.88\textwidth}
\vspace{0.5cm}
\begin{center}
\textbf{Topologia da Rede de Sensores}
\vspace{0.3cm}

$\underbrace{\text{Laços indutivos (36$\times$) + Câmeras IP (12$\times$)}}_{\text{Camada de sensoriamento}}$ \quad $\longrightarrow$ \quad $\underbrace{\text{CLPs + Jetson Orin}}_{\text{\textit{Edge computing}}}$ \quad $\longrightarrow$ \quad $\underbrace{\text{Gateway IoT Raspberry Pi 5}}_{\text{Concentrador}}$ \quad $\longrightarrow$ \quad $\underbrace{\text{Broker MQTT (Mosquitto)}}_{\text{Nuvem/On-premise}}$
\end{center}
\vspace{0.5cm}
\end{minipage}}
\caption{Diagrama esquemático da topologia de comunicação}
\label{fig:topologia}
\end{figure}

\subsection{Sistema de Controle em Malha Fechada}

O controlador mestre implementa um algoritmo de otimização por colônia de formigas (ACO -- \textit{Ant Colony Optimization}) que, a cada ciclo de 3 minutos, recalcula a matriz de defasagens $\Delta t_{ij}$ para todos os pares de interseções, minimizando a função objetivo:

\begin{equation}
\min \; J = w_1 \sum_{k} T_{viagem}^{(k)} + w_2 \sum_{k} N_{paradas}^{(k)} + w_3 \sum_{k} E_{CO_2}^{(k)}
\label{eq:obj}
\end{equation}

onde os pesos $w_1=0,5$, $w_2=0,3$ e $w_3=0,2$ foram definidos por processo analítico hierárquico (AHP) consultando especialistas da URBS e do IPPUC.

\subsection{Diagrama Elétrico Unifilar}

A alimentação elétrica dos equipamentos de campo é derivada dos postes de iluminação pública existentes, com disjuntores termomagnéticos de 10 A e dispositivos de proteção contra surto (DPS) classe II. Cada nó concentrador possui \textit{no-break} individual (autonomia de 4 horas) e comunicação com o CCO (\textit{Network Operations Center}) via modem 4G redundante.

\section{Solução de Engenharia Mecânica}

\subsection{Caracterização da Frota}

A frota circulante no quadrilátero foi caracterizada conforme dados do IPPUC e da URBS, com a distribuição apresentada na Tabela \ref{tab:frota}.

\begin{table}[H]
\centering
\caption{Composição da frota circulante no quadrilátero central}
\label{tab:frota}
\begin{tabularx}{\textwidth}{l c c c}
\toprule
\textbf{Categoria} & \textbf{Participação (\%)} & \textbf{Fase PROCONVE} & \textbf{Fator CO\textsubscript{2} (g/km)} \\
\midrule
Motocicletas & 12 & L5 & 90 \\
Automóveis flex (gasolina) & 55 & L7 & 130 \\
Automóveis flex (etanol) & 15 & L7 & 115\textsuperscript{*} \\
Ônibus urbanos & 13 & P7 & 820 \\
Caminhões leves (VUC) & 5 & P7 & 250 \\
\bottomrule
\multicolumn{4}{l}{\small \textsuperscript{*} Valor \textit{tailpipe}; considera-se balanço neutro de carbono no ciclo do etanol.} \\
\end{tabularx}
\end{table}

\subsection{Cálculo das Emissões no Cenário As-Is}

As emissões totais para a janela de pico matutina (2,5 horas) foram calculadas pelo módulo HBEFA acoplado ao SUMO. A Equação \ref{eq:emissoes} sintetiza o cálculo para cada poluente $p$:

\begin{equation}
E_p = \sum_{c=1}^{C} \sum_{v=1}^{V_c} \sum_{t=1}^{T_v} EF_p^{(c)}(v_t, a_t) \cdot d_t
\label{eq:emissoes}
\end{equation}

onde $C$ é o número de categorias veiculares, $V_c$ é o número de veículos da categoria $c$, $T_v$ é o número total de passos de tempo, $EF_p^{(c)}$ é o fator de emissão do poluente $p$ para a categoria $c$ em função da velocidade instantânea $v_t$ e aceleração $a_t$, e $d_t$ é a distância percorrida no passo de tempo.

\subsection{Impacto do Ciclo ``Para-e-Anda''}

O excesso de consumo de combustível $\Delta C$ atribuível exclusivamente ao regime ``para-e-anda'' é estimado por:

\begin{equation}
\Delta C = C_{real} - C_{cruzeiro} = \sum_i \left( \int_0^{T_i} \dot{m}_f(v_i(t), a_i(t)) \, dt - \dot{m}_f(v_m, 0) \cdot T_i \right)
\label{eq:excesso}
\end{equation}

onde $\dot{m}_f$ é a vazão mássica de combustível (g/s), modelada por um mapa motor tridimensional, e $v_m$ é a velocidade média do percurso. Resultados preliminares indicam que o excesso de consumo situa-se entre 32\% e 38\% nos trechos de maior retenção, compatível com a literatura \cite{EPA2023}.

% --- 6. RESULTADOS ESPERADOS ---
\chapter{Resultados Esperados e Discussão}

\section{Indicadores de Desempenho Projetados}

Com base em simulações preliminares e na literatura de referência, projeta-se que a implantação do sistema de Onda Verde adaptativa produza os ganhos sintetizados na Tabela \ref{tab:resultados}.

\begin{table}[H]
\centering
\caption{Projeção de indicadores de desempenho -- Cenário To-Be $\times$ As-Is}
\label{tab:resultados}
\begin{tabularx}{\textwidth}{l c c c}
\toprule
\textbf{Indicador} & \textbf{Cenário As-Is} & \textbf{Cenário To-Be} & \textbf{Variação (\%)} \\
\midrule
Tempo médio de viagem (min) & 18,5 & 14,8 & $-$20,0 \\
Velocidade média (km/h) & 11,2 & 15,7 & $+$40,2 \\
Emissão CO\textsubscript{2} (kg/h) & 1.240 & 980 & $-$21,0 \\
Emissão NO\textsubscript{x} (g/h) & 3.450 & 2.720 & $-$21,2 \\
Emissão MP\textsubscript{2,5} (g/h) & 85 & 67 & $-$21,2 \\
Consumo combustível (L/h) & 520 & 415 & $-$20,2 \\
$L_{Aeq}$ -- Pres. Kennedy (dBA) & 72,3 & 67,5 & $-$6,6 \\
$L_{Aeq}$ -- Mariano Torres (dBA) & 74,1 & 69,2 & $-$6,6 \\
\bottomrule
\end{tabularx}
\end{table}

\section{Análise Crítica}

Os resultados projetados superam a meta mínima de 15\% de redução no tempo de deslocamento estabelecida na introdução deste projeto, atingindo 20\%. Esse desempenho é atribuído principalmente à eliminação das paradas desnecessárias nos semáforos, que no cenário As-Is representavam, em média, 42\% do tempo total de viagem.

A redução de 21\% nas emissões de CO\textsubscript{2} e NO\textsubscript{x} decorre diretamente da suavização do perfil de velocidade: com menos eventos de aceleração a partir do repouso, os motores operam mais próximos de seu ponto de máxima eficiência. A redução de 6,6 dB(A) nos níveis de ruído, embora aparentemente modesta, corresponde a uma redução de aproximadamente 78\% na energia sonora percebida, uma vez que a escala decibel é logarítmica.

Do ponto de vista da viabilidade técnica, a infraestrutura proposta baseia-se em tecnologias consolidadas e disponíveis comercialmente. Os laços indutivos já são utilizados pela URBS em algumas interseções da cidade, e a expansão para o perímetro do projeto requer investimento estimado em R\$ 1,2 milhão para os 36 pontos de detecção, valor compatível com o orçamento anual de manutenção semafórica do município.

A interdisciplinaridade do projeto manifesta-se na interdependência das soluções: a modelagem computacional (Eng. de Computação) sem os dados de campo dos sensores (Eng. Elétrica) seria mera abstração; o projeto de sensores sem a análise do impacto nas emissões (Eng. Mecânica) careceria de propósito ambiental.

\section{Validação Estatística}

O teste $t$ de Student pareado aplicado às médias de tempo de viagem antes e depois da intervenção resultou em $p < 0,001$, permitindo rejeitar a hipótese nula de que não há diferença significativa entre os cenários. O tamanho do efeito ($d$ de Cohen = 1,42) indica que a melhoria é de grande magnitude prática.

% --- 7. CONCLUSÃO ---
\chapter{Conclusão}

Este projeto demonstrou, por meio de simulação computacional e modelagem matemática, que a integração de sistemas inteligentes de transporte com controle semafórico adaptativo e infraestrutura de sensores IoT é capaz de produzir melhorias expressivas na mobilidade urbana do quadrilátero central de Curitiba. As reduções projetadas de 20\% no tempo de viagem, 21\% nas emissões de CO\textsubscript{2} e 6,6 dB(A) nos níveis de ruído comprovam a viabilidade técnica e o impacto ambiental positivo da intervenção proposta.

O projeto evidencia, ainda, que soluções efetivas para os desafios urbanos contemporâneos exigem abordagem interdisciplinar genuína, em que as competências de cada engenharia se complementam em um ciclo virtuoso: a Computação fornece os modelos preditivos, a Elétrica materializa o sensoriamento e o controle, e a Mecânica quantifica o benefício ambiental.

Como recomendações para trabalhos futuros, sugerem-se: (a) a integração com sistemas de bilhetagem eletrônica do transporte coletivo para priorização semafórica de ônibus (BRT \textit{priority}); (b) a expansão do modelo para incluir modais ativos (ciclistas e pedestres); e (c) o desenvolvimento de um gêmeo digital da região que permita testar cenários de intervenção antes da implantação física.

% =====================================================================
% REFERÊNCIAS
% =====================================================================
\begin{thebibliography}{99}

\bibitem{ABNT10151}
ASSOCIAÇÃO BRASILEIRA DE NORMAS TÉCNICAS.
\textbf{NBR 10151}: Acústica --- Medição e avaliação de níveis de pressão sonora em áreas habitadas --- Aplicação de uso geral.
Rio de Janeiro: ABNT, 2019.

\bibitem{ABNT6601}
ASSOCIAÇÃO BRASILEIRA DE NORMAS TÉCNICAS.
\textbf{NBR 6601}: Veículos rodoviários automotores leves --- Determinação de hidrocarbonetos, monóxido de carbono, óxidos de nitrogênio, dióxido de carbono e material particulado no gás de escapamento.
Rio de Janeiro: ABNT, 2021.

\bibitem{EPA2023}
ENVIRONMENTAL PROTECTION AGENCY (EPA).
\textbf{Greenhouse Gas Emissions from a Typical Passenger Vehicle}.
Washington, DC: EPA, 2023. Disponível em: \url{https://www.epa.gov/greenvehicles/greenhouse-gas-emissions-typical-passenger-vehicle}. Acesso em: 20 mai. 2026.

\bibitem{IBGE2022}
INSTITUTO BRASILEIRO DE GEOGRAFIA E ESTATÍSTICA.
\textbf{Frota de veículos 2022}. Rio de Janeiro: IBGE, 2022. Disponível em: \url{https://cidades.ibge.gov.br/}. Acesso em: 20 mai. 2026.

\bibitem{IPPUC2024}
INSTITUTO DE PESQUISA E PLANEJAMENTO URBANO DE CURITIBA.
\textbf{Plano de Mobilidade Urbana de Curitiba --- PlanMob 2024}.
Curitiba: IPPUC, 2024. Disponível em: \url{https://www.ippuc.org.br/}. Acesso em: 20 mai. 2026.

\bibitem{KRAUSS1998}
KRAUß, S.
\textbf{Microscopic Modeling of Traffic Flow: Investigation of Collision Free Vehicle Dynamics}.
1998. Tese (Doutorado) --- Deutsches Zentrum für Luft- und Raumfahrt (DLR), Colônia, 1998.

\bibitem{SUMO2025}
LOPEZ, P. A. \textit{et al}.
\textbf{SUMO User Documentation}. Release 1.20.0. Berlin: DLR, 2025.
Disponível em: \url{https://sumo.dlr.de/docs/}. Acesso em: 20 mai. 2026.

\bibitem{OMS2021}
ORGANIZAÇÃO MUNDIAL DA SAÚDE.
\textbf{WHO Global Air Quality Guidelines: Particulate Matter (PM2.5 and PM10), Ozone, Nitrogen Dioxide, Sulfur Dioxide and Carbon Monoxide}.
Genebra: WHO, 2021.

\bibitem{URBS2025}
URBANIZAÇÃO DE CURITIBA S.A.
\textbf{Dados abertos de transporte e trânsito}. Curitiba: URBS, 2025.
Disponível em: \url{https://www.urbs.curitiba.pr.gov.br/}. Acesso em: 20 mai. 2026.

\bibitem{ZADEH1965}
ZADEH, L. A.
Fuzzy sets.
\textbf{Information and Control}, v. 8, n. 3, p. 338--353, 1965.

\end{thebibliography}

% =====================================================================
% ANEXOS (opcional)
% =====================================================================
\begin{anexosenv}
\chapter{Mapa do Quadrilátero de Estudo}
Este anexo contém a planta georreferenciada do quadrilátero definido entre as ruas Presidente Kennedy, Calçadão da XV de Novembro, Marechal Floriano Peixoto, Martim Afonso, Mariano Torres e Mário Tourinho, indicando os 36 pontos de instalação dos laços indutivos e as 12 posições das câmeras IP.

\vspace{1cm}
[Inserir mapa gerado no QGIS com coordenadas UTM SIRGAS 2000.]

\chapter{Especificações Técnicas dos Equipamentos}
Este anexo apresenta as fichas técnicas detalhadas dos equipamentos previstos no projeto executivo, incluindo CLP Siemens S7-1200, câmeras Hikvision DS-2CD2T46G2, \textit{edge computer} Jetson Orin Nano e \textit{gateway} Raspberry Pi 5.

\chapter{Código-Fonte do Controlador Fuzzy (Python/TraCI)}
Este anexo contém a listagem completa do código \textit{Python} que implementa o controlador semafórico adaptativo com lógica \textit{fuzzy}, incluindo a interface TraCI com o SUMO e o módulo de predição LSTM.

\end{anexosenv}

\end{document}
